% Part: first-order-logic
% Chapter: natural-deduction
% Section: quantifier-rules

\documentclass[../../../include/open-logic-section]{subfiles}

\begin{document}

\olfileid[fr]{fol}{ntd}{qrl}

\olsection{Règles pour les quantificateurs}

\subsection{Règles pour $\lforall$}

\begin{defish}
\AxiomC{$!A(a)$}
\RightLabel{\Intro{\lforall}}
\UnaryInfC{$\lforall[x][\Atom{!A}{x}]$}
\DisplayProof
\hfill
\AxiomC{$\lforall[x][\Atom{!A}{x}]$}
\RightLabel{\Elim{\lforall}}
\UnaryInfC{$!A(t)$}
\DisplayProof
\end{defish}

Dans les règles pour~$\lforall$, $t$ est un terme clos, c'est-à-dire
un terme sans variable, et $a$~est !!a{constant} qui ne figure ni
dans la conclusion~$\lforall[x][!A(x)]$, ni dans une hypothèse
!!{undischarged} de la !!{derivation} dont la conclusion est la
prémisse~$!A(a)$. On appelle $a$ l'\emph{eigenvariable} de
l'inférence \Intro{\lforall}.\footnote{On emploie le terme
« eigenvariable » pour des raisons historiques, bien que $a$ soit
une constante dans la règle ci-dessus.}

\subsection{Règles pour $\lexists$}

\begin{defish}
\AxiomC{$\Atom{!A}{t}$}
\RightLabel{\Intro{\lexists}}
\UnaryInfC{$\lexists[x][\Atom{!A}{x}]$}
\DisplayProof
\hfill
\AxiomC{$\lexists[x][\Atom{!A}{x}]$}
\AxiomC{[$\Atom{!A}{a}$]$^n$}
\DeduceC{$!C$}
\DischargeRule{\Elim{\lexists}}{n}
\BinaryInfC{$!C$}
\DisplayProof
\end{defish}

De même, $t$ est un terme clos et $a$ est !!a{constant} qui ne
figure ni dans la prémisse $\lexists[x][!A(x)]$, ni dans la
conclusion~$!C$, ni dans une hypothèse !!{undischarged} des
!!{derivation}s dont les conclusions sont les deux prémisses,
à l'exception des hypothèses $!A(a)$.
On appelle $a$ l'\emph{eigenvariable} de l'inférence
\Elim{\lexists}.

Les restrictions que l'on vient d'imposer aux occurrences de la
constante~$a$, pour chacune des règles \Intro{\lforall} et
\Elim{\lexists}, constituent la \emph{condition sur l'eigenvariable}.

\begin{explain}
Rappelons la convention suivante : lorsque $!A$ est !!a{formula}
dans laquelle la !!{variable}~$x$ est libre, on l'indique en
écrivant~$!A(x)$. Dans ce même contexte, $!A(t)$ abrège alors
$\Subst{!A}{t}{x}$. On pourrait donc aussi écrire la règle
$\Intro\lexists$ sous la forme suivante :
\begin{prooftree}
\AxiomC{$\Subst{!A}{t}{x}$}
\RightLabel{\Intro{\lexists}}
\UnaryInfC{$\lexists[x][!A]$}
\end{prooftree}
Remarquons que $t$ peut déjà figurer dans~$!A$; par exemple,
$!A$~pourrait être~$\Atom{\Obj P}{t,x}$. Ainsi, déduire
$\lexists[x][\Atom{\Obj P}{t,x}]$ de~$\Atom{\Obj P}{t,t}$ constitue
une application correcte de~$\Intro\lexists$ : on peut « remplacer »
une ou plusieurs occurrences de~$t$ dans la prémisse, sans
nécessairement les remplacer toutes, par la !!{variable} liée~$x$.
En revanche, les conditions sur l'eigenvariable dans
$\Intro\lforall$ et~$\Elim\lexists$ imposent que la
!!{constant}~$a$ ne figure pas dans~$!A$. On ne peut donc pas
déduire correctement $\lforall[x][\Atom{\Obj P}{a,x}]$ de
$\Atom{\Obj P}{a,a}$ à l'aide de~$\Intro\lforall$.
\end{explain}

\begin{explain}
Dans \Intro{\lexists} et \Elim{\lforall}, aucune autre restriction
ne porte sur le terme clos~$t$. En revanche, dans les règles
\Elim{\lexists} et \Intro{\lforall}, la condition sur
l'eigenvariable impose les restrictions énoncées plus haut aux
occurrences de la !!{constant}~$a$ dans la conclusion et les
hypothèses !!{undischarged}s. Cette condition est nécessaire
à la correction du système : elle garantit que tout !!{sentence}
dérivé est une conséquence des hypothèses !!{undischarged}s.
Sans cette condition, la dérivation suivante serait permise :
\begin{prooftree}
  \AxiomC{$\lexists[x][!A(x)]$}
  \AxiomC{$\Discharge{!A(a)}{1}$}
  \RightLabel{*\Intro{\lforall}}
  \UnaryInfC{$\lforall[x][!A(x)]$}
  \RightLabel{\Elim{\lexists}}
  \BinaryInfC{$\lforall[x][!A(x)]$}
\end{prooftree}
Or, en général, $\lexists[x][!A(x)] \Entails/ \lforall[x][!A(x)]$.

Puisque les règles pour les quantificateurs n'autorisent à
substituer aux !!{variable}s que des termes clos (en particulier
des constantes), toute !!{formula} dérivable à partir d'un ensemble
d'!!{sentence}s est elle-même !!a{sentence}.
\end{explain}

\begin{editorial}
Le résumé de la source interdit à tort toute occurrence de
l'eigenvariable dans les prémisses. On renvoie ici aux conditions
précises données avec chaque règle : $a$ peut notamment figurer
dans $!A(a)$. L'exception concernant les hypothèses $!A(a)$ dans
l'élimination existentielle est conservée telle que la source
l'énonce; elle ne doit pas disparaître dans le rappel ultérieur.
Le terme~$t$
reste clos, même lorsque la source le dit sans restriction. La
condition est donc rappelée par un renvoi à sa formulation complète.
Enfin, la non-conséquence affichée vaut en général, non
pour toute instance de~$!A$, et la dernière explication porte sur
toutes les règles quantificatrices, pas seulement celles d'élimination.
\end{editorial}

\end{document}
